\chapterimage{figs/chapterhead.png}
\chapter{Primeiros passos na interface gráfica}
\label{chap:primeiros_passos_interface_grafica}

\section{Introdução}
\label{sec:cap3_introducao}

Uma vez que a aplicação gráfica do \textit{GenESyS} esteja em funcionamento, o usuário entra, de fato, no simulador. É a partir desse momento que o sistema deixa de ser apenas um projeto compilado e passa a ser um ambiente de trabalho. Este capítulo tem justamente essa função: transformar a primeira execução da GUI em um primeiro contato orientado com a interface.

A proposta aqui não é ainda modelar em detalhe nem executar exemplos completos. Antes disso, o leitor precisa reconhecer a lógica da janela principal, entender quais áreas da interface são mais importantes, perceber onde estão os comandos essenciais e construir uma noção suficientemente clara do espaço de trabalho que o \textit{GenESyS} oferece. Em outras palavras, este capítulo ensina o usuário a se localizar dentro da aplicação.

Essa etapa é mais importante do que pode parecer. Uma interface rica, com múltiplos menus, abas, árvores, painéis e comandos de simulação, pode gerar a impressão de excesso de complexidade quando o usuário tenta compreendê-la de uma só vez. A melhor estratégia é outra: reconhecer primeiro a estrutura global da aplicação, depois identificar seus grupos funcionais e, só então, começar a trabalhar com modelos concretos. É exatamente esse percurso que será seguido aqui.

\section{A janela principal como ambiente de trabalho}
\label{sec:cap3_janela_principal_ambiente}

A interface gráfica do \textit{GenESyS} foi organizada em torno de uma janela principal que centraliza as operações de modelagem, inspeção, simulação, depuração e relatório. Isso significa que o usuário não está diante de uma aplicação com um único painel ou com um conjunto mínimo de comandos isolados. Ao contrário, a GUI foi concebida como um ambiente de trabalho completo, no qual diferentes tarefas podem ser realizadas por meio de diferentes visões do mesmo modelo.

Do ponto de vista do uso, essa janela principal deve ser compreendida como o espaço onde se articulam cinco dimensões fundamentais do trabalho com o simulador:

\begin{enumerate}
    \item a gestão do modelo;
    \item a edição gráfica;
    \item a parametrização e inspeção;
    \item a execução da simulação;
    \item a observação dos resultados e do comportamento do sistema.
\end{enumerate}

Essa organização é especialmente importante porque o usuário não trabalha com o modelo apenas antes da execução nem apenas depois dela. O modelo é construído, ajustado, inspecionado e reavaliado continuamente. A GUI foi estruturada para sustentar exatamente esse ciclo.

\section{Abertura da aplicação e primeira leitura visual}
\label{sec:cap3_abertura_e_leitura_visual}

Ao iniciar a aplicação gráfica, o usuário deve resistir à tentação de clicar imediatamente em todos os menus e comandos disponíveis. O melhor caminho é uma leitura visual orientada da janela principal. Essa leitura pode ser feita em etapas muito simples.

Primeiro, observe a barra de menus e note que ela organiza o trabalho em categorias funcionais. Em geral, estão presentes ações de modelo, edição, visualização, simulação, ferramentas e informações auxiliares. Segundo, observe a área central da interface, que é o espaço principal de representação gráfica do modelo. Terceiro, note a existência de abas, painéis laterais ou inferiores e árvores de navegação, que complementam a visão central. Quarto, observe que a aplicação não foi desenhada apenas para editar um diagrama, mas para acompanhar a vida inteira do modelo, inclusive sua simulação e sua inspeção.

Essa primeira leitura já é suficiente para uma conclusão importante: a GUI do \textit{GenESyS} não é um simples editor gráfico. Ela é uma interface de trabalho integrada.

\section{Grupos funcionais da interface}
\label{sec:cap3_grupos_funcionais_interface}

A forma mais produtiva de compreender a GUI é dividi-la em grupos funcionais. Em vez de decorar comandos isolados, o usuário deve perguntar: que tipo de trabalho cada parte da interface organiza?

\subsection{Comandos de modelo}
\label{subsec:cap3_comandos_modelo}

O primeiro grupo funcional é o dos comandos de modelo. São eles que permitem iniciar um novo trabalho, abrir um modelo já salvo, salvar o modelo atual, fechá-lo, verificar sua consistência e obter informações sobre ele. Esses comandos estruturam o ciclo mais básico de uso da ferramenta.

Na prática, isso significa que o usuário deve reconhecer logo no início onde estão as ações responsáveis por:
\begin{itemize}
    \item criar um novo modelo;
    \item abrir um modelo existente;
    \item salvar o trabalho;
    \item fechar o modelo atual;
    \item verificar o modelo antes de simular.
\end{itemize}

Esses comandos serão usados repetidamente ao longo de toda a primeira parte do livro.

\subsection{Comandos de edição}
\label{subsec:cap3_comandos_edicao}

O segundo grupo é o dos comandos de edição. Aqui entram ações como desfazer, refazer, copiar, colar, excluir, agrupar, desagrupar e, em alguns casos, comandos de organização e alinhamento. Esse conjunto revela que a interface foi pensada para edição efetiva do modelo, e não apenas para observação passiva.

Mesmo que o usuário ainda não esteja construindo modelos novos, é importante perceber desde já que a GUI suporta modificação de elementos gráficos e estruturais. Isso será útil quando os modelos de exemplo começarem a ser alterados, mesmo que de forma leve.

\subsection{Comandos de visualização}
\label{subsec:cap3_comandos_visualizacao}

Outro grupo importante é o dos comandos de visualização. Eles incluem ajustes como zoom, grade, guias e outros recursos que facilitam a organização visual do espaço gráfico. Em um primeiro contato, o usuário não precisa explorar todos esses recursos em detalhe, mas deve reconhecê-los como ferramentas de apoio à leitura e à edição do modelo.

Em especial, o zoom já costuma ser útil logo no início, porque modelos diferentes podem exigir diferentes escalas de observação.

\subsection{Comandos de simulação}
\label{subsec:cap3_comandos_simulacao}

Um grupo central da GUI é formado pelos comandos ligados à execução da simulação. Entre eles estão, tipicamente, iniciar, pausar, retomar, parar, avançar passo a passo e configurar a simulação. Esses comandos mostram que a aplicação não apenas organiza o modelo como estrutura, mas também acompanha sua execução.

Neste capítulo, o leitor deve apenas reconhecer a existência e a localização desses comandos. Seu uso efetivo será aprofundado quando começarmos a trabalhar com modelos concretos. Ainda assim, é importante fixar a ideia de que a mesma GUI usada para abrir e editar o modelo também será usada para conduzir sua simulação.

\subsection{Comandos de inspeção e apoio}
\label{subsec:cap3_comandos_inspecao_apoio}

Além dos grupos anteriores, a interface também oferece mecanismos de inspeção e apoio: painéis de propriedades, árvores de componentes e \textit{data definitions}, áreas de depuração, abas de relatórios e outras visões auxiliares. A presença desses recursos reforça um ponto essencial: o \textit{GenESyS} não foi desenhado apenas para desenhar modelos, mas para estudá-los em execução.

\begin{table}[htbp]
    \centering
    \caption{Grupos funcionais principais da GUI do GenESyS.}
    \label{tab:cap3_grupos_funcionais}
    \small
    \begin{tabular}{|p{4.0cm}|p{7.4cm}|}
        \hline
        \textbf{Grupo funcional} & \textbf{Papel no trabalho do usuário} \\
        \hline
        Modelo & Criar, abrir, salvar, fechar, verificar e organizar o trabalho com modelos \\
        \hline
        Edição & Alterar e reorganizar elementos do modelo \\
        \hline
        Visualização & Ajustar escala, guias e organização visual da área gráfica \\
        \hline
        Simulação & Iniciar, pausar, retomar, interromper e acompanhar a execução \\
        \hline
        Inspeção e apoio & Examinar propriedades, elementos do modelo, resultados e informações auxiliares \\
        \hline
    \end{tabular}
\end{table}

\section{A área gráfica do modelo}
\label{sec:cap3_area_grafica_modelo}

A área gráfica é o centro visual do trabalho com o \textit{GenESyS}. É nela que o modelo aparece de forma diagramática e é nela que a maior parte da atenção do usuário tende a se concentrar durante a leitura e a edição dos modelos. Em um primeiro contato, essa área deve ser compreendida como o espaço principal onde o fluxo do modelo se torna visível.

Isso é importante porque o usuário precisa começar a ler o simulador pela forma como ele mostra o modelo, e não apenas pela forma como o programa organiza seus menus. A área gráfica permite visualizar:
\begin{itemize}
    \item os componentes presentes no modelo;
    \item as conexões entre esses componentes;
    \item a lógica geral do fluxo modelado;
    \item a distribuição visual dos elementos dentro do espaço de trabalho.
\end{itemize}

Mesmo antes de abrir um modelo concreto, o usuário já deve perceber essa área como o núcleo da experiência visual do simulador.

\subsection{Área gráfica não é apenas desenho}
\label{subsec:cap3_area_grafica_nao_apenas_desenho}

Um ponto importante é que a área gráfica não deve ser confundida com um simples editor de desenho. O que se organiza ali não são apenas formas geométricas, mas a representação do modelo que será posteriormente simulado. Em outras palavras, o espaço gráfico é também espaço semântico. Cada elemento visível ali tende a corresponder a uma entidade de modelagem relevante.

Essa observação ajuda o usuário a olhar a interface da forma correta. Em vez de pensar que está apenas manipulando objetos visuais, ele deve entender que está interagindo com a representação operacional do modelo.

\section{Árvores, painéis e editores auxiliares}
\label{sec:cap3_arvores_paineis_editores}

Além da área gráfica central, a GUI do \textit{GenESyS} oferece outros elementos igualmente importantes para o trabalho cotidiano. Esses elementos complementam a leitura visual do modelo e tornam possível inspecioná-lo de forma mais rica.

\subsection{Árvores de navegação}
\label{subsec:cap3_arvores_navegacao}

As árvores de navegação ajudam o usuário a localizar elementos do sistema e do modelo. Dependendo do contexto, elas podem mostrar:
\begin{itemize}
    \item \textit{plugins};
    \item componentes do modelo;
    \item \textit{data definitions};
    \item outras estruturas auxiliares associadas ao trabalho corrente.
\end{itemize}

Em um primeiro contato, essas árvores devem ser vistas como mecanismos de orientação. Elas ajudam a responder à pergunta: que elementos estão disponíveis ou presentes no modelo com o qual estou trabalhando?

\subsection{Editor de propriedades}
\label{subsec:cap3_editor_propriedades}

O editor de propriedades é um dos instrumentos mais importantes da GUI. É por meio dele que o usuário deixa de apenas observar a existência de um elemento e passa a lidar com seus atributos, parâmetros e opções de configuração. Em outras palavras, é nesse ponto que o modelo começa a se tornar efetivamente editável.

Na prática, isso significa que a seleção de um elemento na área gráfica ou em uma árvore de navegação tende a estar associada à exibição de propriedades correspondentes. Ao longo dos próximos capítulos, esse mecanismo aparecerá continuamente, porque quase toda modificação relevante de um modelo passa por algum tipo de parametrização.

\subsection{Abas e visões complementares}
\label{subsec:cap3_abas_visoes_complementares}

A presença de abas centrais e de painéis organizados por abas mostra que a GUI foi construída para sustentar várias perspectivas sobre o mesmo modelo. Algumas dessas perspectivas são visuais, outras são textuais, outras são voltadas à execução, à depuração ou à apresentação de resultados. O usuário não precisa dominar tudo isso agora, mas deve reter a ideia de que o modelo não é observado apenas por uma única superfície visual.

Essa multiplicidade de visões será especialmente importante quando, mais adiante, introduzirmos a aba \textit{Simulang} como representação textual complementar ao modelo gráfico.

\section{Fluxo mínimo de orientação do usuário na GUI}
\label{sec:cap3_fluxo_minimo_orientacao}

Para que a familiarização com a GUI seja produtiva, convém adotar um pequeno roteiro de orientação. Em vez de explorar aleatoriamente a janela principal, o usuário pode seguir esta sequência:

\begin{enumerate}
    \item localizar os comandos de modelo;
    \item identificar a área gráfica central;
    \item localizar árvores e painéis auxiliares;
    \item localizar o editor de propriedades;
    \item reconhecer os comandos básicos de simulação;
    \item observar a existência de abas complementares.
\end{enumerate}

Esse roteiro tem uma vantagem importante: ele evita que a riqueza da interface seja percebida como desorganização. Ao seguir uma ordem simples, o usuário transforma a janela principal em algo inteligível e progressivamente explorável.

\section{O que não é necessário dominar agora}
\label{sec:cap3_o_que_nao_dominar_agora}

Um primeiro contato bem-sucedido com a GUI não exige dominar toda a aplicação. Não é necessário, por exemplo:
\begin{itemize}
    \item usar todos os comandos de edição;
    \item compreender todos os painéis de depuração;
    \item interpretar todas as árvores da interface;
    \item conhecer em profundidade todos os tipos de elementos disponíveis;
    \item explorar ainda a totalidade das linguagens e representações do modelo.
\end{itemize}

Pelo contrário, tentar dominar tudo imediatamente tende a produzir ruído. A aprendizagem mais eficiente ocorre por camadas. Primeiro, o usuário entende a estrutura geral da interface. Depois, passa a trabalhar com modelos concretos. Em seguida, começa a relacionar a GUI com execução, resultados e linguagem textual. Esse é o percurso adotado neste manual.

\section{Boas práticas no primeiro contato com a interface}
\label{sec:cap3_boas_praticas_primeiro_contato}

Há algumas práticas simples que tornam a experiência inicial com a GUI muito mais proveitosa.

A primeira é trabalhar sempre com um objetivo concreto. Em vez de ``explorar a interface'', procure localizar algo específico: o comando de abrir modelo, a área gráfica, o painel de propriedades ou os comandos de simulação.

A segunda é não confundir quantidade de recursos com dificuldade de uso. A interface do \textit{GenESyS} é ampla porque organiza várias fases do trabalho com modelos. Isso não significa que todos os recursos precisem ser usados simultaneamente.

A terceira é observar a interface como ambiente integrado. Menus, área gráfica, árvores, propriedades, simulação e relatórios não são partes independentes. Eles formam um único fluxo de trabalho.

A quarta é aceitar que algumas partes da GUI só passarão a fazer sentido quando um modelo estiver efetivamente carregado. É por isso que o próximo capítulo entrará diretamente nos modelos de exemplo.

\section{Considerações Finais}
\label{sec:cap3_consideracoes_finais}

Este capítulo apresentou a interface gráfica do \textit{GenESyS} como o principal ambiente de trabalho do usuário. Em vez de tratar a GUI como conjunto disperso de comandos, o texto procurou mostrar sua lógica de organização: a janela principal reúne recursos para gerir modelos, editá-los, parametrizá-los, simulá-los e observá-los sob diferentes perspectivas. Essa compreensão é essencial porque evita que o usuário reduza a aplicação a um simples editor gráfico ou, no extremo oposto, a perceba como uma interface excessivamente complexa e desordenada.

Com essa orientação estabelecida, o próximo passo natural é começar a trabalhar sobre modelos concretos. É nesse momento que a interface gráfica realmente ganha sentido operacional. Por isso, o capítulo seguinte passará a tratar dos modelos de exemplo já disponíveis no projeto e de seu uso inicial na GUI, permitindo ao leitor abrir, inspecionar, executar e modificar casos reais de simulação.

Teste seu entendimento desse conteúdo respondendo as seguintes perguntas:

\begin{enumerate}
    \item Por que a GUI do \textit{GenESyS} deve ser entendida como ambiente integrado de trabalho, e não apenas como editor gráfico?

    \item Quais são os grupos funcionais mais importantes da interface em um primeiro contato?

    \item Qual é a função da área gráfica do modelo dentro da experiência de uso do simulador?

    \item Por que não é necessário dominar imediatamente todos os menus, árvores e painéis da aplicação?
\end{enumerate}